\chapter{Design Thinking, Personas e Mapas de Empatia}
\label{cap_design_ux}

A Engenharia de Usabilidade e o \textit{Design Thinking} foram aplicados para garantir que o SIGEA-GTT não apenas cumprisse os requisitos técnicos de software, mas também ressoasse profundamente com as necessidades humanas, cognitivas e pedagógicas de seus usuários reais no ecossistema da Universidade Federal de Sergipe.

\section{Concepção Centrada no Usuário}
\label{sec_design_thinking}

No contexto da educação em saúde e da auditoria clínica, os usuários operam sob alta carga cognitiva, lidando simultaneamente com terminologia médica complexa, protocolos assistenciais rigorosos e a responsabilidade ética pela vida humana. Para mapear com profundidade esse ecossistema, foram desenvolvidas três personas arquetípicas representando os principais grupos de interesse do sistema: a discente de enfermagem em formação, a docente da área da saúde com experiência em gestão da qualidade hospitalar e o docente de computação focado no rigor arquitetural e metodológico.

\section{Especificação das Personas}
\label{sec_personas}

\subsection{Persona 1: Mariana Resende (Discente de Enfermagem)}
\begin{itemize}
    \item \textbf{Dados Demográficos}: 22 anos, solteira, discente do 7º período do Curso de Graduação em Enfermagem da UFS (Campus São Cristóvão / HU-UFS).
    \item \textbf{Contexto e Rotina}: Divide seu tempo entre aulas teóricas no campus universitário e estágios supervisionados nas enfermarias de clínica médica e cirúrgica do Hospital Universitário. Enfrenta rotinas extenuantes de plantão, anotações de enfermagem e administração de medicamentos.
    \item \textbf{Objetivos}: Compreender na prática como os erros assistenciais ocorrem para não repeti-los em sua futura carreira profissional; dominar o raciocínio de causa e efeito e as ferramentas de melhoria contínua; sentir-se confiante na auditoria de prontuários sem receio de punição.
    \item \textbf{Frustrações}: Aulas de segurança do paciente excessivamente teóricas e desprovidas de casos práticos; impossibilidade de manusear prontuários reais de pacientes devido ao sigilo e à LGPD; medo constante de cometer falhas durante os estágios hospitalares.
\end{itemize}

\subsection{Persona 2: Prof.ª Dr.ª Waleska Patrício (Docente da Saúde / Gestora da Qualidade)}
\begin{itemize}
    \item \textbf{Dados Demográficos}: 44 anos, enfermeira com Doutorado em Ciências da Saúde, docente do Departamento de Enfermagem da UFS e colaboradora em comissões de qualidade e segurança do paciente no HU-UFS/EBSERH.
    \item \textbf{Contexto e Rotina}: Ministra disciplinas voltadas à gestão do cuidado, lidera projetos de extensão hospitalar e pesquisa indicadores epidemiológicos de infecções e eventos adversos.
    \item \textbf{Objetivos}: Capacitar turmas numerosas de estudantes na metodologia internacional IHI-GTT em ambiente seguro; dispor de uma ferramenta informatizada para acompanhar o desempenho formativo individual de cada aluno; gerar estatísticas reais de aprendizagem e de incidência de EAs simulados.
    \item \textbf{Frustrações}: Uso de planilhas eletrônicas obsoletas e descentralizadas para tabulação de casos de aula; sobrecarga na correção manual de estudos de caso em papel; carência de sistemas computacionais brasileiros adaptados à realidade do SUS e do hospital universitário.
\end{itemize}

\subsection{Persona 3: Prof. Dr. Gilton Ferreira (Docente e Pesquisador DCOMP)}
\begin{itemize}
    \item \textbf{Dados Demográficos}: 48 anos, Doutor em Ciência da Computação, docente titular do Departamento de Computação (DCOMP/UFS) e orientador acadêmico.
    \item \textbf{Contexto e Rotina}: Leciona disciplinas de Engenharia de Software, Arquitetura de Sistemas e Banco de Dados; orienta projetos de graduação e pós-graduação voltados à informática em saúde e sistemas de missão crítica.
    \item \textbf{Objetivos}: Garantir que o sistema de software seja projetado sob rigorosos padrões da engenharia (arquitetura desacoplada, alta coesão, 100\% de cobertura de testes, segurança e persistência confiável); fomentar a interdisciplinaridade concreta entre a Computação e a Saúde na UFS.
    \item \textbf{Frustrações}: Softwares acadêmicos desenvolvidos como protótipos frágeis (\textit{throwaway prototypes}) que são abandonados após a defesa do TCC; sistemas sem testes automatizados ou com arquiteturas monolíticas de difícil manutenção.
\end{itemize}

\section{Mapas de Empatia}
\label{sec_mapas_empatia}

Os Mapas de Empatia estruturados para cada uma das personas encontram-se consolidados nos Quadros \ref{qua_empatia_mariana}, \ref{qua_empatia_waleska} e \ref{qua_empatia_gilton}.

\begin{quadro}[htb]
\centering
\small
\caption{Mapa de Empatia - Persona 1: Mariana Resende (Discente de Enfermagem).}
\label{qua_empatia_mariana}
\begin{tabularx}{\textwidth}{p{3.8cm} X}
\toprule
\textbf{Dimensão Empática} & \textbf{Síntese Observada no Ecossistema Universitário / Hospitalar} \\
\midrule
\textbf{O que pensa e sente?} & Sente receio de errar uma dosagem ou não perceber a piora clínica de um paciente; deseja ser uma enfermeira segura, crítica e humanizada; pensa que a teoria da sala de aula é distante da correria dos plantões hospitalares. \\
\midrule
\textbf{O que escuta?} & Ouve preceptores alertando sobre os riscos de punição e processos éticos do COREN; ouve colegas comentando sobre incidentes ocorridos nas enfermarias que foram abafados por medo; escuta professores enfatizando a importância das metas internacionais da OMS. \\
\midrule
\textbf{O que vê?} & Vê prontuários volumosos com anotações manuais por vezes ilegíveis; vê equipes multidisciplinares sobrecarregadas e com pouco tempo para auditoria; vê colegas com dúvidas conceituais na identificação de eventos adversos. \\
\midrule
\textbf{O que fala e faz?} & Tenta revisar duas vezes as medicações antes de administrá-las; debate com colegas de estágio sobre condutas assistenciais; manifesta interesse por atividades práticas que simulem situações críticas da vida real. \\
\midrule
\textbf{Dores (Frustrações)} & Medo paralisante da culpabilização individual; insegurança ao diferenciar uma reação adversa esperada de uma falha assistencial evitável; falta de acesso a prontuários completos para estudo. \\
\midrule
\textbf{Necessidades (Ganhos)} & Uma plataforma de simulação realista onde possa treinar a busca de gatilhos sem colocar vidas em risco; feedback pedagógico detalhado do professor explicando seus acertos e equívocos; domínio prático de ferramentas como Ishikawa e GUT. \\
\bottomrule
\end{tabularx}
\fonte{Elaborado pelo autor (2026).}
\end{quadro}

\begin{quadro}[htb]
\centering
\small
\caption{Mapa de Empatia - Persona 2: Prof.ª Dr.ª Waleska Patrício (Docente da Saúde).}
\label{qua_empatia_waleska}
\begin{tabularx}{\textwidth}{p{3.8cm} X}
\toprule
\textbf{Dimensão Empática} & \textbf{Síntese Observada no Ecossistema Universitário / Hospitalar} \\
\midrule
\textbf{O que pensa e sente?} & Acredita fervorosamente que a mudança cultural nos hospitais começa na graduação; sente-se frustrada quando os alunos chegam ao internato sem dominar conceitos básicos de segurança do paciente; preocupa-se com o cumprimento das normas da LGPD. \\
\midrule
\textbf{O que escuta?} & Escuta a direção do HU cobrando indicadores de qualidade e acreditação hospitalar; ouve relatos de subnotificação endêmica nos setores de internação; ouve dos discentes que as aulas práticas com prontuários reais são restritas. \\
\midrule
\textbf{O que vê?} & Vê o IHI-GTT sendo aplicado em hospitais de excelência internacional com resultados extraordinários; vê a carência de um software nacional gratuito voltado especificamente ao ensino universitário; observa a dificuldade dos alunos em estruturar planos de ação 5W2H e ciclos PDCA. \\
\midrule
\textbf{O que fala e faz?} & Defende que ``errar em ambiente simulado é aprender, errar no paciente é imperdoável''; elabora estudos de caso clínicos ricos em detalhes; busca constantemente parcerias interdisciplinares com a Computação. \\
\midrule
\textbf{Dores (Frustrações)} & Tempo excessivo despendido na gestão manual de turmas e correções de atividades em papel; dificuldade de mensurar o avanço coletivo da turma em indicadores como $TI_{1000}$ e $TI_{100}$. \\
\midrule
\textbf{Necessidades (Ganhos)} & Um ambiente web completo onde possa cadastrar prontuários simulados complexos; painéis de controle da turma com métricas automáticas do IHI; módulo de correção formativa com histórico e parecer pedagógico estruturado. \\
\bottomrule
\end{tabularx}
\fonte{Elaborado pelo autor (2026).}
\end{quadro}

\begin{quadro}[htb]
\centering
\small
\caption{Mapa de Empatia - Persona 3: Prof. Dr. Gilton Ferreira (Docente e Pesquisador DCOMP).}
\label{qua_empatia_gilton}
\begin{tabularx}{\textwidth}{p{3.8cm} X}
\toprule
\textbf{Dimensão Empática} & \textbf{Síntese Observada no Ecossistema Universitário / Hospitalar} \\
\midrule
\textbf{O que pensa e sente?} & Valoriza a Engenharia de Software de alto nível aplicada a problemas sociais concretos; pensa que sistemas na área da saúde exigem tolerância zero a falhas arquiteturais e inconsistências de banco de dados; preza pela formação científica rigorosa do formando. \\
\midrule
\textbf{O que escuta?} & Escuta a academia mundial demandando soluções de software de código limpo (\textit{Clean Code}), segurança rigorosa e conformidade com privacidade; escuta órgãos de fomento valorizando a interdisciplinaridade entre Computação e Saúde. \\
\midrule
\textbf{O que vê?} & Vê projetos acadêmicos que falham por negligenciarem testes de software e documentação técnica; vê a evolução das tecnologias web modernas (Java 21 LTS, Spring Boot 3.3, Angular 18 e PostgreSQL) como a base ideal para sistemas robustos. \\
\midrule
\textbf{O que fala e faz?} & Exige qualidade máxima nos critérios de aceite de código (100\% de cobertura de testes automatizados e compilação sem alertas); incentiva a publicação de artigos científicos e registro de propriedade intelectual; orienta com foco em padrões arquiteturais canônicos. \\
\midrule
\textbf{Dores (Frustrações)} & Projetos de graduação que não alcançam maturidade de produção; documentações técnicas incompletas ou desconformes com os padrões do DCOMP-UFS. \\
\midrule
\textbf{Necessidades (Ganhos)} & Uma plataforma de software estado-da-arte, com arquitetura documentada no modelo 4+1, diagramas formais em TikZ, cobertura universal de testes e validação empírica pronta para publicação de impacto científico. \\
\bottomrule
\end{tabularx}
\fonte{Elaborado pelo autor (2026).}
\end{quadro}
